% titel: Blatt 3, A1: Natriumcarbonat aus NaOH und CO2 (limitierendes Reagenz)
% skript: Mortimer Kap. 4.5 (Stöchiometrie, limitierendes Reagenz)

% Teil I: Werkzeuge, kein Lösungsweg.

Werden zwei Edukte im falschen Verhältnis eingesetzt, bleibt eines davon
teilweise übrig: Das zuerst verbrauchte Edukt (das limitierende Reagenz)
bestimmt, wie viel Produkt maximal entstehen kann.

\begin{werkzeugkasten}[Stoffmenge und limitierendes Reagenz]
Aus Masse und molarer Masse die Stoffmenge jedes Edukts berechnen:
\begin{equation}
  n = \frac{m}{M}.
\end{equation}
Für eine ausgeglichene Reaktion
$a\,\mathrm{A} + b\,\mathrm{B} \rightarrow \mathrm{Produkt}$ ist das Edukt
mit dem \emph{kleineren} Quotienten $n/\nu$ ($\nu$ = Stöchiometriefaktor) das
limitierende Reagenz. Aus seiner Stoffmenge folgt über den Faktor $a$ die
maximal mögliche Stoffmenge des Produkts; die Masse erhält man mit
$m = n\cdot M$.

\textit{Skript:} Mortimer Kap. 4 (Reaktionsgleichungen, Stöchiometrie),
Übung 3/1.
\end{werkzeugkasten}

\begin{achtung}[Darauf achten]
\begin{itemize}
  \item Beide Edukt-Stoffmengen berechnen und die Quotienten $n/\nu$
        vergleichen --- nicht aus den \emph{Massen} allein schließen.
  \item Die Reaktionsgleichung muss zuerst sauber ausgeglichen sein; die
        Stöchiometriefaktoren (z.\,B. $2\,\mathrm{NaOH} : 1\,\mathrm{CO_2}$)
        stecken in den Nennern.
  \item Die Produktmenge aus dem limitierenden Edukt rechnen; das
        überschüssige Edukt bleibt als Rest übrig.
  \item Einheiten $(\mathrm{g/mol},\ \mathrm{mol})$ durchgehend mitführen.
\end{itemize}
\end{achtung}
